\secnumberlesssection{ANEXOS}

En los Anexos se incluye todo aquel material complementario que no es parte del contenido de los capítulos de la memoria, pero que permiten a un lector contar con un contenido adjunto relacionado con el tema.

\subsection{ENCUESTA SOBRE TECNOLOGÍAS PARA PERSONAS CIEGAS}
\label{anexo:encuesta}

Este anexo detalla el instrumento mencionado en el Capítulo 1 (Definición del problema), aplicado para validar empíricamente la problemática social antes de proponer una solución.

\begin{table}[h]
    \centering
    \caption{\label{table:fichaencuesta} Ficha técnica de la encuesta.} Fuente: Elaboración propia.
    \begin{tabular}{|p{5cm}|p{9cm}|}
        \hline
        Instrumento & ``Encuesta sobre tecnologías para personas ciegas'' \\
        \hline
        Plataforma & Google Forms \\
        \hline
        Validación previa & Pedro Correa Maturana, director del INLAB (Unidad de Inclusión Laboral) de FUNDALURP (Fundación de Lucha contra la Retinitis Pigmentosa) \\
        \hline
        Canal de distribución & Correo electrónico, a personas asistidas por FUNDALURP o vinculadas a ella \\
        \hline
        Período de aplicación & 26 de febrero de 2026 -- 11 de marzo de 2026 \\
        \hline
        Respuestas obtenidas & 9 (N=9) \\
        \hline
        Encuestados que usan lector de pantalla & 7 de 9 (78\%) \\
        \hline
    \end{tabular}
\end{table}

La encuesta se estructuró en cuatro bloques temáticos:

\begin{itemize}
    \item \textbf{Datos generales y de contacto}: marca temporal de la respuesta, correo electrónico opcional (para diferenciar participantes y evitar respuestas duplicadas), e interés en ser contactado a futuro en el contexto del proyecto.
    \item \textbf{Perfil del encuestado}: género; año de nacimiento (usado para estimar el contexto tecnológico según la época); condición que causa la discapacidad visual (la más severa o antigua, si el encuestado padece más de una); forma en que la condición afecta la visión; nivel de discapacidad visual; grado de manejo computacional actual. Si la discapacidad fue adquirida en una etapa avanzada de la vida (y no de nacimiento), se agregan dos preguntas adicionales: el grado de manejo computacional y el tipo de tecnologías que dominaba antes de la pérdida de visión.
    \item \textbf{Uso de lectores de pantalla}: si el encuestado conoce qué es y qué hace un lector de pantalla; si emplea uno para utilizar su computador; cuáles conoce (aunque no los use); cuál usa como lector principal; y qué impacto tuvo en su vida comenzar a usarlo. Las preguntas de este bloque y del siguiente solo se muestran a quienes reportan emplear un lector de pantalla.
    \item \textbf{Dificultades y funcionalidades}: dificultades típicas al utilizar un lector de pantalla en general; dificultades típicas al navegar sitios web específicamente con uno; selección, de una lista cerrada, de las 3 funcionalidades que le parecerían más útiles; una pregunta abierta opcional para sugerir una funcionalidad adicional no considerada en la lista; y una pregunta abierta opcional de cierre para comentarios sobre la encuesta.
\end{itemize}

\subsubsection{Respuestas completas}

A continuación se listan las 9 respuestas completas de la encuesta. Cada participante se identifica con un código (P1 a P9) en lugar de su correo electrónico, dato personal que se omite por completo. El orden de los participantes corresponde al orden cronológico de respuesta.

Dos preguntas del formulario usan categorías cerradas cuyo texto completo se abrevia en la tabla para no repetirlo en cada fila; su definición completa es la siguiente:

\begin{itemize}
    \item \textbf{Nivel de discapacidad visual}: \textit{Moderada} (dificultad para leer textos estándar y para reconocer rostros a corta distancia, incluso con corrección); \textit{Severa} (no logra leer textos impresos convencionales y requiere apoyos específicos para la orientación y el desplazamiento en espacios públicos); \textit{Ceguera total} (no presenta visión funcional y no puede leer textos sin el uso de apoyos como lectores de pantalla u otros sistemas de accesibilidad).
    \item \textbf{Manejo computacional}: \textit{Básico} (utiliza el computador para navegar por la web, enviar correos y acceder a redes sociales); \textit{Intermedio} (utiliza el computador de forma cotidiana y maneja herramientas de productividad como Word y Excel); \textit{Avanzado} (utiliza software especializado para el trabajo y realiza ajustes en configuraciones del sistema o de las aplicaciones); \textit{Experto} (domina el uso del computador y es capaz de enseñar a otras personas cómo utilizarlo).
\end{itemize}

\begin{longtable}{|p{4cm}|p{10.5cm}|}
    \caption{\label{table:respuestascompletas} Respuestas completas de la encuesta, anonimizadas. Fuente: Elaboración propia.} \\
    \hline
    \textbf{Campo} & \textbf{Respuesta} \\
    \hline
    \endfirsthead
    \hline
    \textbf{Campo} & \textbf{Respuesta} \\
    \hline
    \endhead
    \hline
    \endfoot

    \multicolumn{2}{|l|}{\textbf{Participante P1}} \\
    \hline
    Género & Femenino \\
    \hline
    Año de nacimiento & 1974 \\
    \hline
    Condición & Retinitis pigmentaria \\
    \hline
    Cómo afecta la visión & Ceguera total \\
    \hline
    Nivel de discapacidad visual & Ceguera total \\
    \hline
    Manejo computacional actual & Avanzado \\
    \hline
    Origen de la discapacidad & Adquirida (por enfermedad o accidente) \\
    \hline
    Manejo computacional antes de la pérdida & Experto \\
    \hline
    Tecnologías dominadas antes de la pérdida & Navegar por sitios web, redes sociales (Facebook, Twitter, Instagram), correo electrónico (Gmail, Outlook), tiendas en línea (Falabella, AliExpress), banca en línea, videoconferencias (Zoom, Google Meet), productividad (Word, Excel), manejo de carpetas y archivos del sistema, programas de diseño gráfico, edición audiovisual o arquitectura (Adobe, AutoDesk) \\
    \hline
    ¿Conoce qué es un lector de pantalla? & Sí \\
    \hline
    ¿Emplea un lector de pantalla? & Sí \\
    \hline
    Lectores de pantalla que conoce & JAWS, NVDA, TalkBack, Narrador \\
    \hline
    Lector de pantalla principal & NVDA \\
    \hline
    Impacto de comenzar a usarlo & (sin respuesta) \\
    \hline
    Dificultades al utilizar un lector de pantalla & Incapacidad de interpretar elementos estrictamente visuales como imágenes; incapacidad de leer notificaciones o \textit{pop-ups} independientes de la estructura del sitio o ubicados en otra sección de la página que no está seleccionada en ese momento \\
    \hline
    Dificultades al navegar sitios web & Recuadros de texto a rellenar sin etiqueta o que no informan qué información se debe ingresar; obligación de usar el ratón para interactuar con el sitio; pruebas de verificación humana estrictamente visuales \\
    \hline
    Funcionalidades más útiles seleccionadas & Describir e interpretar imágenes, fotografías y gráficos; solicitar por voz la navegación directa a una sección del sitio (ej. ``Llévame al precio''); lectura de notificaciones o etiquetas dinámicas \\
    \hline
    Funcionalidad adicional sugerida & (sin respuesta) \\
    \hline
    Comentario final & (sin respuesta) \\
    \hline

    \multicolumn{2}{|l|}{\textbf{Participante P2}} \\
    \hline
    Género & Femenino \\
    \hline
    Año de nacimiento & 1975 \\
    \hline
    Condición & Retinitis pigmentaria \\
    \hline
    Cómo afecta la visión & Visión de túnel o periférica reducida; pérdida de visión central; sensibilidad a la luz \\
    \hline
    Nivel de discapacidad visual & Severa \\
    \hline
    Manejo computacional actual & Intermedio \\
    \hline
    Origen de la discapacidad & Adquirida (por enfermedad o accidente) \\
    \hline
    Manejo computacional antes de la pérdida & Intermedio \\
    \hline
    Tecnologías dominadas antes de la pérdida & Navegar por sitios web, redes sociales (Facebook, Twitter, Instagram), correo electrónico (Gmail, Outlook), tiendas en línea (Falabella, AliExpress), banca en línea, videoconferencias (Zoom, Google Meet), productividad (Word, Excel), manejo de carpetas y archivos del sistema \\
    \hline
    ¿Conoce qué es un lector de pantalla? & Sí \\
    \hline
    ¿Emplea un lector de pantalla? & No \\
    \hline

    \multicolumn{2}{|l|}{\textbf{Participante P3}} \\
    \hline
    Género & Masculino \\
    \hline
    Año de nacimiento & 1965 \\
    \hline
    Condición & Degeneración Macular Asociada a la Edad (DMAE) \\
    \hline
    Cómo afecta la visión & Visión borrosa o nublada en general; zonas oscuras o manchas flotantes; sensibilidad a la luz \\
    \hline
    Nivel de discapacidad visual & Severa \\
    \hline
    Manejo computacional actual & Básico \\
    \hline
    Origen de la discapacidad & Adquirida (por enfermedad o accidente) \\
    \hline
    Manejo computacional antes de la pérdida & Intermedio \\
    \hline
    Tecnologías dominadas antes de la pérdida & Navegar por sitios web, redes sociales (Facebook, Twitter, Instagram), correo electrónico (Gmail, Outlook), tiendas en línea (Falabella, AliExpress), banca en línea, manejo de carpetas y archivos del sistema, programas de diseño gráfico, edición audiovisual o arquitectura (Adobe, AutoDesk) \\
    \hline
    ¿Conoce qué es un lector de pantalla? & Sí \\
    \hline
    ¿Emplea un lector de pantalla? & Sí \\
    \hline
    Lectores de pantalla que conoce & TalkBack, Narrador \\
    \hline
    Lector de pantalla principal & VoiceOver \\
    \hline
    Impacto de comenzar a usarlo & ``Fue como volver a ser más autónomo'' \\
    \hline
    Dificultades al utilizar un lector de pantalla & Incapacidad de interpretar elementos estrictamente visuales como imágenes; orden de lectura ilógico que obliga a pasar por información irrelevante antes de llegar a la clave; incapacidad de leer notificaciones o \textit{pop-ups} \\
    \hline
    Dificultades al navegar sitios web & Recuadros de texto sin etiqueta; el sistema entrega información valiosa mediante imágenes sin texto alternativo; pruebas de verificación humana estrictamente visuales \\
    \hline
    Funcionalidades más útiles seleccionadas & Obtener un resumen del contenido principal del sitio; poder hacer preguntas sobre el sitio (ej. ``¿Qué puedo hacer en esta página?''); solicitar por voz la navegación directa a una sección del sitio \\
    \hline
    Funcionalidad adicional sugerida & (sin respuesta) \\
    \hline
    Comentario final & (sin respuesta) \\
    \hline

    \multicolumn{2}{|l|}{\textbf{Participante P4}} \\
    \hline
    Género & Masculino \\
    \hline
    Año de nacimiento & 1972 \\
    \hline
    Condición & Nanoftalmo \\
    \hline
    Cómo afecta la visión & Visión borrosa o nublada en general; baja visión, imposibilidad de identificar figuras, letras o detalles \\
    \hline
    Nivel de discapacidad visual & Severa \\
    \hline
    Manejo computacional actual & Básico \\
    \hline
    Origen de la discapacidad & De nacimiento \\
    \hline
    ¿Conoce qué es un lector de pantalla? & No \\
    \hline
    ¿Emplea un lector de pantalla? & No \\
    \hline

    \multicolumn{2}{|l|}{\textbf{Participante P5}} \\
    \hline
    Género & Masculino \\
    \hline
    Año de nacimiento & 1973 \\
    \hline
    Condición & Opacidad de córnea (síndrome de Peters) \\
    \hline
    Cómo afecta la visión & Visión borrosa o nublada en general; distingue colores y bultos, aunque los colores muy oscuros los percibe como negro (ej. distingue semáforos de noche o con clima nublado) \\
    \hline
    Nivel de discapacidad visual & Severa \\
    \hline
    Manejo computacional actual & Básico \\
    \hline
    Origen de la discapacidad & De nacimiento \\
    \hline
    ¿Conoce qué es un lector de pantalla? & Sí \\
    \hline
    ¿Emplea un lector de pantalla? & Sí \\
    \hline
    Lectores de pantalla que conoce & JAWS, NVDA, TalkBack \\
    \hline
    Lector de pantalla principal & NVDA \\
    \hline
    Impacto de comenzar a usarlo & ``Me siento en igualdad de condiciones ante las personas sin discapacidad visual al usar un pc como ellos'' \\
    \hline
    Dificultades al utilizar un lector de pantalla & Incapacidad de interpretar elementos estrictamente visuales como imágenes; orden de lectura ilógico; incapacidad de leer notificaciones o \textit{pop-ups}; nula reacción ante cambios en la presentación del sitio (secciones que aparecen y desaparecen); respuesta libre: ``los gráficos y cuando no especifica los botones cualquieras éstos sean'' \\
    \hline
    Dificultades al navegar sitios web & El sistema entrega información valiosa mediante imágenes sin texto alternativo; instrucciones estrictamente visuales (ej. ``seleccione la opción verde''); la página cambia su estructura sin avisar; obligación de usar el ratón; pruebas de verificación humana estrictamente visuales; respuesta libre: ``los gráficos y botones sin descripción'' \\
    \hline
    Funcionalidades más útiles seleccionadas & Describir e interpretar imágenes, fotografías y gráficos; solicitar por voz la navegación directa a una sección del sitio; lectura de notificaciones o etiquetas dinámicas \\
    \hline
    Funcionalidad adicional sugerida & ``El significado de los botones y gráficos que muestran los símbolos e imágenes pero para el lector no están descritos en forma escrita, como para que nos diga: botón ingresar, la imagen representa tal o cual figura... cuestiones así por el estilo'' \\
    \hline
    Comentario final & ``Descripción de gestos, textos, gráficos e imágenes'' \\
    \hline

    \multicolumn{2}{|l|}{\textbf{Participante P6}} \\
    \hline
    Género & Masculino \\
    \hline
    Año de nacimiento & 2008 \\
    \hline
    Condición & Neuropatía óptica hereditaria de Leber \\
    \hline
    Cómo afecta la visión & Pérdida de visión central; escasa visión periférica y cambio en la percepción de colores \\
    \hline
    Nivel de discapacidad visual & Severa \\
    \hline
    Manejo computacional actual & Básico \\
    \hline
    Origen de la discapacidad & Adquirida (por enfermedad o accidente) \\
    \hline
    Manejo computacional antes de la pérdida & Básico \\
    \hline
    Tecnologías dominadas antes de la pérdida & Navegar por sitios web, videoconferencias (Zoom, Google Meet) \\
    \hline
    ¿Conoce qué es un lector de pantalla? & No \\
    \hline
    ¿Emplea un lector de pantalla? & Sí \\
    \hline
    Lectores de pantalla que conoce & NVDA, TalkBack \\
    \hline
    Lector de pantalla principal & TalkBack \\
    \hline
    Impacto de comenzar a usarlo & ``Me ayudó a adaptarme a esta nueva normalidad de mejor forma'' \\
    \hline
    Dificultades al utilizar un lector de pantalla & Orden de lectura ilógico que obliga a pasar por información irrelevante antes de llegar a la clave \\
    \hline
    Dificultades al navegar sitios web & La página cambia su estructura u organización sin avisar \\
    \hline
    Funcionalidades más útiles seleccionadas & Obtener un resumen del contenido principal del sitio; describir e interpretar imágenes, fotografías y gráficos; solicitar por voz la navegación directa a una sección del sitio \\
    \hline
    Funcionalidad adicional sugerida & (sin respuesta) \\
    \hline
    Comentario final & (sin respuesta) \\
    \hline

    \multicolumn{2}{|l|}{\textbf{Participante P7}} \\
    \hline
    Género & Otro \\
    \hline
    Año de nacimiento & 1982 \\
    \hline
    Condición & Glaucoma \\
    \hline
    Cómo afecta la visión & (respuesta del encuestado poco clara en el formulario original: ``Lucero visión'') \\
    \hline
    Nivel de discapacidad visual & Ceguera total \\
    \hline
    Manejo computacional actual & Intermedio \\
    \hline
    Origen de la discapacidad & De nacimiento \\
    \hline
    ¿Conoce qué es un lector de pantalla? & Sí \\
    \hline
    ¿Emplea un lector de pantalla? & Sí \\
    \hline
    Lectores de pantalla que conoce & JAWS, NVDA, VoiceOver, TalkBack, Narrador, Orca \\
    \hline
    Lector de pantalla principal & VoiceOver \\
    \hline
    Impacto de comenzar a usarlo & ``Es mi forma y manera de comunicarme con otros y ser productivo laboralmente'' \\
    \hline
    Dificultades al utilizar un lector de pantalla & Incapacidad de interpretar elementos estrictamente visuales como imágenes; orden de lectura ilógico; incapacidad de leer notificaciones o \textit{pop-ups}; nula reacción ante cambios en la presentación del sitio \\
    \hline
    Dificultades al navegar sitios web & Recuadros de texto sin etiqueta; imágenes sin texto alternativo; navegar textos extensos sin estructura; instrucciones estrictamente visuales; la página cambia su estructura sin avisar; obligación de usar el ratón; pruebas de verificación humana estrictamente visuales \\
    \hline
    Funcionalidades más útiles seleccionadas & Obtener un resumen del contenido principal del sitio; poder hacer preguntas sobre el sitio; describir e interpretar imágenes, fotografías y gráficos \\
    \hline
    Funcionalidad adicional sugerida & (sin respuesta) \\
    \hline
    Comentario final & (sin respuesta) \\
    \hline

    \multicolumn{2}{|l|}{\textbf{Participante P8}} \\
    \hline
    Género & Femenino \\
    \hline
    Año de nacimiento & 1979 \\
    \hline
    Condición & Glaucoma \\
    \hline
    Cómo afecta la visión & Ceguera total \\
    \hline
    Nivel de discapacidad visual & Ceguera total \\
    \hline
    Manejo computacional actual & Básico \\
    \hline
    Origen de la discapacidad & Adquirida (por enfermedad o accidente) \\
    \hline
    Manejo computacional antes de la pérdida & Intermedio \\
    \hline
    Tecnologías dominadas antes de la pérdida & Navegar por sitios web \\
    \hline
    ¿Conoce qué es un lector de pantalla? & Sí \\
    \hline
    ¿Emplea un lector de pantalla? & Sí \\
    \hline
    Lectores de pantalla que conoce & JAWS, NVDA, VoiceOver, TalkBack \\
    \hline
    Lector de pantalla principal & NVDA \\
    \hline
    Impacto de comenzar a usarlo & (sin respuesta) \\
    \hline
    Dificultades al utilizar un lector de pantalla & Incapacidad de interpretar elementos estrictamente visuales como imágenes; orden de lectura ilógico; incapacidad de leer notificaciones o \textit{pop-ups} \\
    \hline
    Dificultades al navegar sitios web & Recuadros de texto sin etiqueta; imágenes sin texto alternativo; navegar textos extensos sin estructura; instrucciones estrictamente visuales; pruebas de verificación humana estrictamente visuales \\
    \hline
    Funcionalidades más útiles seleccionadas & Obtener un resumen del contenido principal del sitio; describir e interpretar imágenes, fotografías y gráficos; solicitar por voz la navegación directa a una sección del sitio \\
    \hline
    Funcionalidad adicional sugerida & (sin respuesta) \\
    \hline
    Comentario final & (sin respuesta) \\
    \hline

    \multicolumn{2}{|l|}{\textbf{Participante P9}} \\
    \hline
    Género & Femenino \\
    \hline
    Año de nacimiento & 1967 \\
    \hline
    Condición & Retinitis pigmentaria \\
    \hline
    Cómo afecta la visión & Visión de túnel o periférica reducida; zonas oscuras o manchas flotantes \\
    \hline
    Nivel de discapacidad visual & Moderada \\
    \hline
    Manejo computacional actual & Básico \\
    \hline
    Origen de la discapacidad & Adquirida (por enfermedad o accidente) \\
    \hline
    Manejo computacional antes de la pérdida & Básico \\
    \hline
    Tecnologías dominadas antes de la pérdida & Productividad (Word, Excel) \\
    \hline
    ¿Conoce qué es un lector de pantalla? & Sí \\
    \hline
    ¿Emplea un lector de pantalla? & Sí \\
    \hline
    Lectores de pantalla que conoce & NVDA, TalkBack \\
    \hline
    Lector de pantalla principal & NVDA \\
    \hline
    Impacto de comenzar a usarlo & ``Me emocioné mucho porque, cuando escribo, el lector me va leyendo y así estoy segura de lo que voy escribiendo'' \\
    \hline
    Dificultades al utilizar un lector de pantalla & Respuesta libre: ``mi dificultad es que el lector no deja de hablar'' \\
    \hline
    Dificultades al navegar sitios web & Navegar textos extensos sin estructura; instrucciones estrictamente visuales \\
    \hline
    Funcionalidades más útiles seleccionadas & Poder hacer preguntas sobre el sitio; describir e interpretar imágenes, fotografías y gráficos; solicitar por voz la navegación directa a una sección del sitio \\
    \hline
    Funcionalidad adicional sugerida & (sin respuesta) \\
    \hline
    Comentario final & ``Me parece una iniciativa muy buena que se sigan buscando mejoras para el lector de pantalla'' \\
    \hline
\end{longtable}

\subsection{CALIBRACIÓN DEL JUEZ LLM: HISTORIAL DE VERSIONES DEL \textit{PROMPT}}
\label{anexo:promptsjuez}

Este anexo detalla el proceso mencionado en el Capítulo 3 (sección \ref{sec:benchmark}, ``Selección del juez''): la calibración del \textit{prompt} del juez LLM contra el conjunto de 40 respuestas evaluadas por un humano, a lo largo de 23 versiones (v1 a v23). Para el candidato gemma3-27b se incluye además el desglose de aciertos por nivel de la rúbrica (L0/L1/L2/L3, sobre un total de 10 casos por nivel).

\begin{longtable}{|p{1.5cm}|p{5cm}|p{2.8cm}|p{2.8cm}|}
    \caption{\label{table:historialprompts} Resultados de \textit{exact match} por versión de \textit{prompt} y modelo juez. Fuente: Elaboración propia.} \\
    \hline
    \textbf{Versión} & \textbf{gemma3-27b (L0/L1/L2/L3)} & \textbf{gemma-3-12b} & \textbf{MPrometheus-14B} \\
    \hline
    \endfirsthead
    \hline
    \textbf{Versión} & \textbf{gemma3-27b (L0/L1/L2/L3)} & \textbf{gemma-3-12b} & \textbf{MPrometheus-14B} \\
    \hline
    \endhead
    \hline
    \endfoot
    v1 & sin métricas formales & -- & -- \\
    \hline
    v2 & sin métricas formales & -- & -- \\
    \hline
    v3 & sin métricas formales & -- & -- \\
    \hline
    v4 & sin métricas formales & -- & -- \\
    \hline
    v5 & 70,0\% (7/3/9/9) & 60,0\% & 45,0\% \\
    \hline
    v6 & 70,0\% (4/6/8/10) & 57,5\% & 52,5\% \\
    \hline
    v7 & 67,5\% (4/6/7/10) & 52,5\% & 45,0\% \\
    \hline
    v8 & 62,5\% (3/5/7/10) & 55,0\% & 50,0\% \\
    \hline
    v9 & \textbf{72,5\% (5/6/8/10)} & 60,0\% & 52,5\% \\
    \hline
    v10 & 72,5\% (4/6/9/10) & 55,0\% & 42,5\% \\
    \hline
    v11 & 70,0\% (5/5/8/10) & 62,5\% & 57,5\% \\
    \hline
    v12 & 70,0\% (4/5/9/10) & 57,5\% & 47,5\% \\
    \hline
    v13 & 70,0\% (5/5/8/10) & -- & -- \\
    \hline
    v14 & 67,5\% (3/6/8/10) & 50,0\% & 45,0\% \\
    \hline
    v15 & 42,5\% (5/8/2/2) & 37,5\% & 32,5\% \\
    \hline
    v16 & 70,0\% (6/4/8/10) & -- & -- \\
    \hline
    v17 & 70,0\% (6/4/8/10) & -- & -- \\
    \hline
    v18 & 70,0\% (6/4/8/10) & -- & -- \\
    \hline
    v19 & 67,5\% (6/3/8/10) & -- & -- \\
    \hline
    v20 & \textbf{72,5\% (6/5/8/10)} & -- & -- \\
    \hline
    v21 & 70,0\% (6/6/6/10) & -- & -- \\
    \hline
    v22 & 67,5\% (6/4/7/10) & -- & -- \\
    \hline
    v23 & 70,0\% (6/6/6/10) & -- & -- \\
    \hline
\end{longtable}

\subsection{DETALLE DE LAS SESIONES DE ENTREVISTAS}
\label{anexo:sesiones}

Este anexo detalla el proceso mencionado en el Capítulo 4 (sección de entrevistas con usuarios): las 4 sesiones realizadas sobre los 3 participantes, con su modalidad, duración y los sitios evaluados en cada una.

\begin{table}[h]
    \centering
    \caption{\label{table:detallesesiones} Detalle de las sesiones de entrevistas.} Fuente: Elaboración propia.
    \begin{tabular}{|p{1.15cm}|p{1.8cm}|p{1.9cm}|p{1.6cm}|p{4.75cm}|p{1.4cm}|}
        \hline
        \textbf{Sesión} & \textbf{Usuario} & \textbf{Modalidad} & \textbf{Duración} & \textbf{Sitios evaluados} & \textbf{N° tareas} \\
        \hline
        0 & Usuario 0 & Texto & 30 min & USM (educacional), Falabella (\textit{e-commerce}) & 4 \\
        \hline
        1 & Usuario 1 & Texto & 1 hora & YouTube (redes sociales), portal docente de una institución educativa (educacional) & 9 \\
        \hline
        2 & Usuario 2 & Texto & 15 min & Portal docente y portal de estudiante de instituciones educativas (educacional), SII (estatal) & 5 \\
        \hline
        3 & Usuario 2 & Voz & 15 min & SII (estatal) & 1 \\
        \hline
    \end{tabular}
\end{table}

\subsection{REGISTRO DE LA PRUEBA DE RECONOCIMIENTO DE VOZ (STT)}
\label{anexo:stt}

Este anexo detalla el proceso mencionado en el Capítulo 4 (sección de métricas técnicas, exactitud del reconocimiento de voz): las 10 frases leídas en voz alta por cada una de las 5 personas que participaron en la prueba, la transcripción literal devuelta por el motor de STT de la extensión, y la cantidad de palabras erróneas de cada una respecto del total de palabras de la frase original.

\begin{longtable}{|p{0.5cm}|p{5.7cm}|p{5.7cm}|p{1.5cm}|}
    \caption{\label{table:transcripcionesstt} Registro completo de la prueba de reconocimiento de voz: frase original, transcripción obtenida y palabras erróneas por frase, para cada una de las 5 personas.} \\
    \hline
    \textbf{\#} & \textbf{Frase original} & \textbf{Transcripción obtenida} & \textbf{Palabras erróneas} \\
    \hline
    \endfirsthead
    \hline
    \textbf{\#} & \textbf{Frase original} & \textbf{Transcripción obtenida} & \textbf{Palabras erróneas} \\
    \hline
    \endhead
    \hline
    \endfoot

    \multicolumn{4}{|l|}{\textbf{Persona 1}} \\
    \hline
    1 & ¿Dónde puedo encontrar los canales de contacto oficial de la institución? & Dónde puedo encontrar los canales de contacto oficial de la institución & 0 / 11 \\
    \hline
    2 & Indícame, de haber, dónde debo acceder para ingresar a mi cuenta. & Indícame a ver dónde debo acceder para ingresar a mi cuenta & 2 / 11 \\
    \hline
    3 & ¿En qué parte exacta de la pantalla debo ubicarme para revisar el carrito de compras? & En qué parte exacta la pantalla debe ubicarme para revisar el cargo de compras & 3 / 15 \\
    \hline
    4 & ¿Qué forma tienen los elementos interactuables ubicados debajo de cada publicación, y qué propósito tienen? & Qué forma tienen los elementos intelectual de ubicados debajo de cada publicación Y qué propósito tiene & 3 / 15 \\
    \hline
    5 & ¿Cuál es la temática de este sitio web, y qué tipo de contenido se visualiza en él? & Cuál es la temática de este sitio web Y qué tipo de contenido se visualiza en él & 0 / 17 \\
    \hline
    6 & Basado en la forma estructural del sitio, ¿cómo se interactúa físicamente con la página para seguir consumiendo contenido? & basado en la forma estructural del sitio Cómo se interactúa físicamente con la página para seguir consumiendo contenido & 0 / 18 \\
    \hline
    7 & ¿Dónde se encuentra el menú principal de navegación y qué opciones o portales clave destaca para estudiantes o postulantes? & Dónde se encuentra el menú principal de navegación y Qué opciones o portales clave destaca para sus estudiantes de postulante & 3 / 19 \\
    \hline
    8 & ¿Cuál es el propósito principal de esta página web y qué tipo de información o servicios públicos se ofrecen en ella? & Cuál es el propósito principal de esta página web Y qué tipo de información o servicios públicos se ofrecen en ella & 0 / 21 \\
    \hline
    9 & ¿De qué trata esta página web y cuáles son los bloques o secciones que componen su estructura principal de arriba hacia abajo? & De qué trata esta página web y cuáles son los bloqueos secciones que componen su estructura principal de arriba hacia abajo & 2 / 22 \\
    \hline
    10 & ¿Dónde se ubican espacialmente las herramientas de accesibilidad del sitio (si las hay) y qué funciones de ajuste visual o lectura ofrecen? & Dónde se ubican espacialmente la herramienta de accesibilidad del sitio si las hay y qué funciones de ajuste visual la lectura ofrecen & 3 / 22 \\
    \hline

    \multicolumn{4}{|l|}{\textbf{Persona 2}} \\
    \hline
    1 & ¿Dónde puedo encontrar los canales de contacto oficial de la institución? & Dónde puedo encontrar los canales de contacto oficial de la institución & 0 / 11 \\
    \hline
    2 & Indícame, de haber, dónde debo acceder para ingresar a mi cuenta. & Indícame a ver dónde puedo acceder para ingresar a mi cuenta & 3 / 11 \\
    \hline
    3 & ¿En qué parte exacta de la pantalla debo ubicarme para revisar el carrito de compras? & En qué parte exacta de la pantalla debo ubicarme para revisar el carrito de compras & 0 / 15 \\
    \hline
    4 & ¿Qué forma tienen los elementos interactuables ubicados debajo de cada publicación, y qué propósito tienen? & Qué forma tiene los elementos interactuables ubicados debajo de cada publicación Y qué propósito tienen & 1 / 15 \\
    \hline
    5 & ¿Cuál es la temática de este sitio web, y qué tipo de contenido se visualiza en él? & Cuál es la temática de este sitio web Y qué tipo de contenido se visualiza en él & 0 / 17 \\
    \hline
    6 & Basado en la forma estructural del sitio, ¿cómo se interactúa físicamente con la página para seguir consumiendo contenido? & basado en la fórmula estructural del sirio Cómo se interactúa físicamente con la página para seguir consumiendo contenido & 2 / 18 \\
    \hline
    7 & ¿Dónde se encuentra el menú principal de navegación y qué opciones o portales clave destaca para estudiantes o postulantes? & Dónde se encuentra el menú principal de navegación y Qué opciones o portales clave destaca para estudiantes o postulantes & 0 / 19 \\
    \hline
    8 & ¿Cuál es el propósito principal de esta página web y qué tipo de información o servicios públicos se ofrecen en ella? & Cuál es el propósito principal de esta página web Y qué tipo de información o servicios públicos se ofrecen en ella & 0 / 21 \\
    \hline
    9 & ¿De qué trata esta página web y cuáles son los bloques o secciones que componen su estructura principal de arriba hacia abajo? & De qué trata esta página web y cuáles son los bloques o secciones que componen su estructura principal de arriba hacia abajo & 0 / 22 \\
    \hline
    10 & ¿Dónde se ubican espacialmente las herramientas de accesibilidad del sitio (si las hay) y qué funciones de ajuste visual o lectura ofrecen? & Dónde se ubican especialmente las herramientas de accesibilidad del sitio y Qué funciones de ajuste visual o lectura ofrecen & 4 / 22 \\
    \hline

    \multicolumn{4}{|l|}{\textbf{Persona 3}} \\
    \hline
    1 & ¿Dónde puedo encontrar los canales de contacto oficial de la institución? & Dónde puedo encontrar los canales de contacto oficial de la institución & 0 / 11 \\
    \hline
    2 & Indícame, de haber, dónde debo acceder para ingresar a mi cuenta. & Indícame de haber dónde debo acceder para ingresar a mi cuenta & 0 / 11 \\
    \hline
    3 & ¿En qué parte exacta de la pantalla debo ubicarme para revisar el carrito de compras? & En qué parte exacta de la pantalla debo ubicarme para revisar el carrito de compras & 0 / 15 \\
    \hline
    4 & ¿Qué forma tienen los elementos interactuables ubicados debajo de cada publicación, y qué propósito tienen? & Qué forma tienen los elementos interactuables ubicados debajo de cada publicación Y qué propósito tienen & 0 / 15 \\
    \hline
    5 & ¿Cuál es la temática de este sitio web, y qué tipo de contenido se visualiza en él? & Cuál es la temática de este sitio web Y qué tipo de contenido se visualiza en él & 0 / 17 \\
    \hline
    6 & Basado en la forma estructural del sitio, ¿cómo se interactúa físicamente con la página para seguir consumiendo contenido? & basado en la forma estructural del sitio Cómo se interactúa físicamente con la página para seguir consumiendo contenido & 0 / 18 \\
    \hline
    7 & ¿Dónde se encuentra el menú principal de navegación y qué opciones o portales clave destaca para estudiantes o postulantes? & Dónde se encuentra el menú principal de navegación y Qué opciones soportales clave destaca para estudiantes o postulantes & 2 / 19 \\
    \hline
    8 & ¿Cuál es el propósito principal de esta página web y qué tipo de información o servicios públicos se ofrecen en ella? & Cuál es el propósito principal de esta página web Y qué tipo de información o servicios públicos se ofrecen en ella & 0 / 21 \\
    \hline
    9 & ¿De qué trata esta página web y cuáles son los bloques o secciones que componen su estructura principal de arriba hacia abajo? & De qué trata esta página web y cuáles son los bloques o secciones que componen su estructura principal de arriba hacia abajo & 0 / 22 \\
    \hline
    10 & ¿Dónde se ubican espacialmente las herramientas de accesibilidad del sitio (si las hay) y qué funciones de ajuste visual o lectura ofrecen? & Dónde se ubican especialmente las herramientas de accesibilidad del sitio si las hay y qué funciones de ajuste visual o lectura ofrecen & 1 / 22 \\
    \hline

    \multicolumn{4}{|l|}{\textbf{Persona 4}} \\
    \hline
    1 & ¿Dónde puedo encontrar los canales de contacto oficial de la institución? & Dónde puedo encontrar los canales de contacto oficial de la institución & 0 / 11 \\
    \hline
    2 & Indícame, de haber, dónde debo acceder para ingresar a mi cuenta. & Indícame de a ver dónde debo acceder para ingresar a mi cuenta & 2 / 11 \\
    \hline
    3 & ¿En qué parte exacta de la pantalla debo ubicarme para revisar el carrito de compras? & En qué parte exacta de la pantalla debo ubicarme para revisar el carrito de compras & 0 / 15 \\
    \hline
    4 & ¿Qué forma tienen los elementos interactuables ubicados debajo de cada publicación, y qué propósito tienen? & Qué forma tienen los elementos interactuables ubicados debajo de cada publicación Y qué propósito tienen & 0 / 15 \\
    \hline
    5 & ¿Cuál es la temática de este sitio web, y qué tipo de contenido se visualiza en él? & Cuál es la temática de este sitio web Y qué tipo de contenido se visualiza en él & 0 / 17 \\
    \hline
    6 & Basado en la forma estructural del sitio, ¿cómo se interactúa físicamente con la página para seguir consumiendo contenido? & basado en la forma estructural del sitio Cómo se interactúa físicamente con la página para seguir consumiendo contenido & 0 / 18 \\
    \hline
    7 & ¿Dónde se encuentra el menú principal de navegación y qué opciones o portales clave destaca para estudiantes o postulantes? & Dónde se encuentra el menú principal de navegación y Qué opciones o portales clave destaca para estudiantes o postulantes & 0 / 19 \\
    \hline
    8 & ¿Cuál es el propósito principal de esta página web y qué tipo de información o servicios públicos se ofrecen en ella? & Cuál es el propósito principal de esta página web Y qué tipo de información o servicios públicos se ofrecen en ella & 0 / 21 \\
    \hline
    9 & ¿De qué trata esta página web y cuáles son los bloques o secciones que componen su estructura principal de arriba hacia abajo? & De qué trata esta página web y cuáles son los bloques o secciones que componen su estructura principal de arriba hacia abajo & 0 / 22 \\
    \hline
    10 & ¿Dónde se ubican espacialmente las herramientas de accesibilidad del sitio (si las hay) y qué funciones de ajuste visual o lectura ofrecen? & Dónde se ubican espacialmente las herramientas de accesibilidad del sitio si las hay y qué funciones de ajuste visual o lectura ofrecen & 0 / 22 \\
    \hline

    \multicolumn{4}{|l|}{\textbf{Persona 5}} \\
    \hline
    1 & ¿Dónde puedo encontrar los canales de contacto oficial de la institución? & Dónde puedo encontrar los canales de contacto oficial de la institución & 0 / 11 \\
    \hline
    2 & Indícame, de haber, dónde debo acceder para ingresar a mi cuenta. & Indícame a ver dónde debo acceder para ingresar a mi cuenta & 2 / 11 \\
    \hline
    3 & ¿En qué parte exacta de la pantalla debo ubicarme para revisar el carrito de compras? & En qué parte salta de la pantalla de ubicarme para revisar el carrito de compras & 2 / 15 \\
    \hline
    4 & ¿Qué forma tienen los elementos interactuables ubicados debajo de cada publicación, y qué propósito tienen? & Qué forma tiene los elementos intelectuales ubicados debajo de cada publicación Y qué propósito tienen & 2 / 15 \\
    \hline
    5 & ¿Cuál es la temática de este sitio web, y qué tipo de contenido se visualiza en él? & Cuál es la temática de este sitio web Y qué tipo de contenido se visualiza en él & 0 / 17 \\
    \hline
    6 & Basado en la forma estructural del sitio, ¿cómo se interactúa físicamente con la página para seguir consumiendo contenido? & basado en la forma estructural del sitio Cómo se interactúa físicamente con la página para seguir consumiendo contenido & 0 / 18 \\
    \hline
    7 & ¿Dónde se encuentra el menú principal de navegación y qué opciones o portales clave destaca para estudiantes o postulantes? & Dónde se encuentra el menú principal de navegación y Qué opciones o portales clave destaca para estudiantes y postulantes & 1 / 19 \\
    \hline
    8 & ¿Cuál es el propósito principal de esta página web y qué tipo de información o servicios públicos se ofrecen en ella? & Cuál es el propósito principal de esta página web Y qué tipo de información o servicios públicos se ofrecen en ella & 0 / 21 \\
    \hline
    9 & ¿De qué trata esta página web y cuáles son los bloques o secciones que componen su estructura principal de arriba hacia abajo? & De qué trata esta página web y cuáles son los bloques o secciones que componen su estructura principal de arriba hacia abajo & 0 / 22 \\
    \hline
    10 & ¿Dónde se ubican espacialmente las herramientas de accesibilidad del sitio (si las hay) y qué funciones de ajuste visual o lectura ofrecen? & Dónde se ubican especialmente las herramientas de accesibilidad del sitio y Qué funciones de ajuste visual o lectura ofrecen & 4 / 22 \\
    \hline
\end{longtable}
